\section{Metodología}

\subsection{Enfoque experimental}

La estrategia metodológica consiste en descomponer el problema general —dotar a un asistente virtual de capacidades de hiperpersonalización— en situaciones atómicas que puedan resolverse de forma aislada mediante técnicas de \textit{prompting} sobre modelos de lenguaje. Con los hallazgos obtenidos en cada situación se avanza, de manera incremental, hacia una solución general. El trabajo es de naturaleza experimental: cada decisión de arquitectura se valida observando el comportamiento del sistema sobre casos de uso concretos antes de incorporarla a la solución.

\subsection{Etapas de trabajo}

El trabajo de investigación se organiza en las siguientes etapas sucesivas, donde cada fase refina la arquitectura de recuperación y generación de respuestas:

\begin{enumerate}
	\item Experimentos de \textit{prompting} sobre problemas de alcance muy focalizado (menos de diez ejemplos), utilizando datos generados ex profeso.
	\item Creación de conjuntos de datos a partir de documentación real.
	\item Generación de conjuntos de datos sintéticos.
	\item Generalización hacia conjuntos de datos de alcance intermedio (aproximadamente un año) y su almacenamiento en bases de datos con capacidad de recuperación contextual.
\end{enumerate}

\subsection{Patrón de desarrollo}

Para cada problemática abordada se siguió un patrón iterativo:

\begin{enumerate}
	\item Se seleccionó una problemática enunciada como un caso de uso imposible de resolver hoy sin hiperpersonalización.
	\item Se describió una estrategia de ataque y se desarrolló un prototipo.
	\item Se registraron los resultados para alimentar un repositorio de \textit{prompts} exitosos.
	\item Se repitió el proceso hasta encontrar una solución satisfactoria.
\end{enumerate}
